\section{Metode}

\begin{figure}[t]
    \centering
    \resizebox{0.78\columnwidth}{!}{%
    \begin{tikzpicture}[
        node distance=0.42cm,
        block/.style={rectangle, draw, rounded corners, fill=blue!7,
                      minimum height=0.52cm, text width=3.05cm,
                      align=center, font=\scriptsize},
        arr/.style={-{Stealth[length=2.5pt]}, thick}
    ]
        \node[block] (data) {Dataset Vincze et al.};
        \node[block, below=of data] (filter) {Filtering nilai inti kosong};
        \node[block, below=of filter] (features) {Hitung log massa, log umur, exposure proxy};
        \node[block, below=of features] (reg) {Regresi + bootstrap + robustness};
        \node[block, below=of reg] (resid) {Risiko naif + residual supresi};
        \node[block, below=of resid] (bio) {Interpretasi biologis siklus sel};
        \draw[arr] (data) -- (filter);
        \draw[arr] (filter) -- (features);
        \draw[arr] (features) -- (reg);
        \draw[arr] (reg) -- (resid);
        \draw[arr] (resid) -- (bio);
    \end{tikzpicture}}
    \caption{Bagan alir metode komputasi. Analisis bergerak dari dataset
    komparatif menuju regresi, uji robustness, model risiko naif, dan
    interpretasi mekanistik.}
    \label{fig:method-flow}
\end{figure}

\subsection{Data dan Filtering}

Data utama adalah \texttt{SupplementaryData.xls} dari Vincze dkk.
\cite{vincze2022}, diunduh dari repositori GitHub penulis dan dikonversi
menjadi CSV. Dataset memuat spesies mamalia, ordo, massa tubuh, adult life
expectancy, CMR, ICM, jumlah individu, jumlah individu mati, jumlah kasus
neoplasia, dan jumlah catatan postmortem. Baris dengan nilai kosong pada
massa tubuh, adult life expectancy, CMR, atau ICM dikeluarkan. Semua analisis
berikutnya menggunakan 172 spesies yang lolos filtering.

Kode Python tersedia pada repositori proyek:
\url{https://github.com/Nayekah/KDS}. Implementasi memakai pustaka standar
Python agar analisis dapat dijalankan tanpa instalasi paket numerik eksternal.
Komponen program terdiri dari pembaca data, fungsi statistik, analisis
komparatif, dan generator tabel/figur untuk laporan.

% Generated by python main.py
\begin{table}[t]
\centering
\caption{Ringkasan empiris dataset mamalia sebagai konteks Peto's Paradox.}
\label{tab:empirical-context}
\begin{tabular}{lr}
\toprule
Besaran & Nilai\\
\midrule
Jumlah spesies & 172\\
Spesies CMR $>$ 0 & 131\\
Slope log massa--log CMR & -0.086\\
$R^2$ log massa--log CMR & 0.053\\
Slope exposure--log CMR & -0.063\\
$R^2$ exposure--log CMR & 0.034\\
\bottomrule
\end{tabular}
\end{table}

\newcommand{\SpeciesCount}{172}
\newcommand{\NonzeroSpeciesCount}{131}
\newcommand{\MassSlope}{-0.086}
\newcommand{\MassRsq}{0.053}
\newcommand{\ExposureSlope}{-0.063}
\newcommand{\ExposureRsq}{0.034}

CMR digunakan sebagai variabel risiko utama karena merepresentasikan
probabilitas kematian terkait kanker pada individu dewasa. Adult life
expectancy dikonversi dari hari menjadi tahun. Massa tubuh digunakan dalam
kilogram. Untuk mengurangi efek skala yang lebar, hubungan statistik
dihitung pada skala logaritmik untuk massa tubuh, adult life expectancy,
exposure proxy, dan CMR nonnol.

\subsection{Exposure Proxy}

Risiko kanker secara teoritis meningkat dengan jumlah sel dan lama paparan
waktu. Karena jumlah sel aktual tidak tersedia untuk semua spesies, massa
tubuh digunakan sebagai proksi jumlah sel. Exposure proxy didefinisikan
sebagai
\begin{equation}
    E_i = M_i \times L_i,
\end{equation}
dengan $M_i$ adalah massa tubuh spesies $i$ dan $L_i$ adalah adult life
expectancy dalam tahun. Proksi ini tidak mengukur jumlah pembelahan aktual,
tetapi cukup untuk menguji prediksi dasar Peto's Paradox: spesies dengan
$E$ besar diperkirakan memiliki risiko kanker lebih tinggi jika tidak ada
kompensasi biologis.

\subsection{Regresi Skala Tubuh dan Risiko}

Analisis pertama menghitung regresi linear sederhana:
\begin{equation}
    \log_{10}(\mathrm{CMR}_i) =
    \alpha + \beta \log_{10}(X_i) + \epsilon_i,
\end{equation}
dengan $X_i$ berupa massa tubuh, adult life expectancy, atau exposure proxy.
Hanya spesies dengan CMR lebih dari nol yang digunakan dalam regresi log.
Nilai slope dan $R^2$ dilaporkan untuk melihat apakah risiko kanker
meningkat searah dengan ukuran tubuh dan umur.

Untuk menguji kestabilan slope, bootstrap dilakukan 1000 kali dengan
resampling spesies ber-CMR nonnol. Dari distribusi slope bootstrap diambil
rerata dan interval kepercayaan 95\%. Selain itu, regresi massa tubuh
diulang setelah menerapkan ambang minimum catatan postmortem 20, 50, dan
100. Uji ini digunakan karena CMR pada spesies dengan sedikit catatan
postmortem lebih rentan terhadap noise pencatatan.

\subsection{Model Risiko Naif dan Residual Supresi}

Untuk menafsirkan deviasi dari prediksi naif, risiko berbasis exposure
didefinisikan sebagai
\begin{equation}
    R_i^\mathrm{naif} = 1 - \exp(-\lambda E_i).
\end{equation}
Parameter $\lambda$ dikalibrasi sehingga median risiko naif mendekati median
CMR nonnol. Residual supresi kemudian dihitung sebagai
\begin{equation}
    S_i = 1 -
    \frac{-\ln(1-\mathrm{CMR}_i)}{\lambda E_i}.
\end{equation}
Nilai $S_i$ tinggi menunjukkan CMR teramati lebih rendah daripada risiko naif
yang diprediksi oleh exposure proxy. Nilai ini bukan pengukuran langsung
DNA repair atau apoptosis, tetapi kandidat sinyal supresi kanker relatif
yang dapat dikaitkan dengan mekanisme kontrol siklus sel.

\subsection{Ringkasan Ordo}

Selain analisis spesies, data diringkas per ordo. Ordo dengan kurang dari
tiga spesies dikeluarkan dari ringkasan agar median tidak terlalu rapuh.
Untuk setiap ordo dihitung median CMR, median massa tubuh, median adult life
expectancy, dan jumlah spesies. Ringkasan ini digunakan untuk melihat
apakah perbedaan risiko lebih tampak pada tingkat taksonomi daripada pada
skala tubuh saja.

\subsection{Kriteria Interpretasi}

Hasil kuantitatif ditafsirkan dalam tiga tingkat. Pertama, slope regresi
digunakan untuk menjawab apakah ukuran tubuh, umur, atau exposure proxy
memiliki hubungan arah yang konsisten dengan CMR. Kedua, bootstrap dan filtering
postmortem digunakan untuk menilai apakah kesimpulan tersebut stabil
terhadap sampling spesies dan kualitas catatan. Ketiga, residual supresi
digunakan sebagai analisis kualitatif untuk memilih spesies kandidat.
Dengan pemisahan ini, paper tidak menyamakan korelasi statistik dengan
mekanisme biologis langsung, tetapi menggunakan hasil komputasi sebagai
dasar formulasi hipotesis biologis yang lebih terarah.
